\section{Introduction}
\label{sec:intro}

Piezoelectric semiconductors, of which n-type zinc oxide is the standard example, join a piezoelectric response to a population of mobile charge carriers, a combination that underlies nanogenerators and piezotronic devices \cite{wang2006nanogenerators} and rests on a well-characterized material \cite{ozgur2005zno}.  Acoustic waves in such a medium are not purely elastic: the piezoelectric field that accompanies a wave acts on the carriers, which redistribute, screen the field and absorb energy from the wave \cite{hutson1962elastic}, and drifting carriers can even amplify ultrasound \cite{white1962amplification}.  The linearized drift--diffusion description of the carriers has since been combined with the elastic and piezoelectric equations for rods \cite{zhang2016extensional} and plates \cite{zhang2025pscplate}, and it is the description adopted here.

When the plate is a few hundred nanometers thick and the frequency is of the order of ten gigahertz, two further effects can no longer be set aside.  The first is flexoelectricity, the polarization produced by a strain gradient rather than by a strain \cite{tagantsev1986piezo,majdoub2008enhanced,yudin2013fundamentals,zubko2013flexoelectric,shen2010theory}, which scales inversely with a length and therefore becomes comparable with piezoelectricity in thin films and near free surfaces and edges; it is already exploited in semiconductors \cite{bhaskar2016flexoelectric}.  The second is the dependence of the elastic energy on the strain gradient, which introduces an internal length into the equations of motion \cite{mindlin1964micro}.  In a semiconductor the two effects are coupled to the carriers, because the bound charge induced by the strain gradient is screened by mobile charge and the redistributed carriers act back on the field.  This closed loop is the content of the flexoelectric semiconductor theory \cite{qu2020mobile,qu2022thermal,qu2022bending}, whose most complete statement is the thermodynamic derivation of Qu and co-workers \cite{qu2025thermodynamic}; a computational treatment of flexoelectric semiconductor junctions has also appeared \cite{hu2025computational}.

Wave propagation in flexoelectric media has been studied for several geometries.  Lamb waves in a plate with flexoelectricity and strain-gradient elasticity \cite{yang2018lamb} and Rayleigh waves on a flexoelectric half-space \cite{yang2020rayleigh} have been computed, and plane-wave reflection and transmission have been obtained for a flexoelectric slab \cite{jiao2018slab}, for a piezoelectric semiconductor slab with carrier and bias effects \cite{jiao2019semiconductor}, and at the free surface of a pre-stressed flexoelectric half-space \cite{mondal2022scattering}.  For flexoelectric and piezoelectric semiconductors, Rayleigh waves \cite{biswas2025rayleigh,gupta2025psvoid} and the exact multi-field solution for a plate \cite{zhang2025pscplate} have been analyzed, and guided waves in plates with strain-gradient elasticity alone are available for the Mindlin Form-II and generalized first strain-gradient theories \cite{vavva2009gradient,sidhardh2019microplate}.  In all of these studies the quantity computed is a dispersion relation or a set of reflection coefficients for the few body waves of a half-space.

The scattering of a guided mode at a free plate edge is, by contrast, a classical problem for purely elastic plates.  The edge couples the propagating Lamb modes to an infinite family of evanescent ones, and the partition of the incident power among the propagating modes was obtained by modal expansion and variational methods \cite{torvik1967reflection,auld1977variational,gregory1983reflection}, compared with measurement \cite{morvan2003lamb} and extended to oblique incidence \cite{santhanam2013reflection}.  The same edge supports a resonance at frequencies of zero group velocity \cite{prada2005laser}, where the dispersion branch turns back on itself.  Lamb modes include backward modes, whose energy flux is opposed to their phase velocity \cite{meitzler1965backward}; where a forward and a backward mode coexist near a zero-group-velocity point \cite{prada2008local}, the free edge converts one into the other and produces negative reflection \cite{germano2012anomalous,gerardin2016negative,veres2016broad}.  The same conversion at thickness steps gives negative refraction \cite{bramhavar2011negative}, in the manner of the metamaterial negative refraction of elastic waves \cite{zhu2014negative}.  Computationally, guided waves in plates are now routinely handled by semi-analytical methods that discretize the thickness and leave the propagation direction analytic \cite{adamou2004spectral,bartoli2006modeling}, and the Legendre--Galerkin construction \cite{shen1994efficient} gives well-conditioned systems for the fourth-order operators that strain-gradient elasticity introduces.

This leaves a gap.  The edge-scattering results above are all for plates without strain-gradient elasticity, flexoelectricity or mobile carriers, whereas the flexoelectric-semiconductor results do not treat a bounded plate: they give dispersion and attenuation of modes, or reflection of body waves at a half-space, and none resolves how an incident guided mode is distributed among the outgoing branches of a free edge.  The two questions are connected, because in a strain-gradient plate the outgoing spectrum is discrete and may contain a backward member, and because the carriers can bias the partition through the attenuation they give the participating modes.  An edge partition is also a different calculation from a dispersion curve: it requires matching a continuum of outgoing modes to the traction-free edge conditions and classifying each of them, and the generalized traction and the independent double traction of a strain-gradient plate must both be imposed.

The present article addresses this gap for a free edge of an n-type ZnO plate described by strain-gradient elasticity, piezoelectricity, flexoelectricity and linearized drift--diffusion, with one additional carrier--strain-gradient coupling that is constrained by a thermodynamic potential (\cref{sec:governing}).  Two questions are answered separately: whether the plate admits a backward branch at the frequencies of interest and how strongly the edge converts the incident mode into it, and whether the carrier--strain-gradient coupling biases that partition.  The main contributions are (i) a formulation of the free-edge problem in which the generalized traction and the double traction are imposed separately; (ii) a spectral Legendre--Galerkin solution validated against an independent $C^{1}$ finite-element discretization, a two-dimensional finite-element solution, exact limits and an energy balance; (iii) the partition of the incident power into a backward branch, with the distinction between a modal share and a reflectance made explicit; and (iv) a quantified, bounded effect of the added coupling on that partition, reported with its own controls so that it can be read separately from the effect of the branch itself.

Throughout, the plate occupies $x<0$ with the reflecting edge at $x=0$, and the $z$ axis is the six-fold axis of ZnO, through the thickness, with the faces at $z=\pm d/2$.  A mode is called \emph{backward} when its energy flux has a negative $x$ component, and \emph{outgoing} when it either decays away from the edge or, if propagating, carries energy away from it.  The formulation and the solution method are given first, followed by the numerical results, and the findings are collected in the conclusions.
